\section{Checks}
\label{sec:checks}
This section gives a structured account of the verifications performed on the chain
exact observable $\to$ field equations $\to$ line-of-sight functionals $\to$ observer bookkeeping $\to$ harmonic machinery.
Each item points to the verification ledger (App.~\ref{app:ledger}), where the code location and the numbers are recorded.
Checks of the assembled $C^{(22)}_\ell$, $C^{(13)}_\ell$ and of the monochromatic evaluation belong to
Secs.~\ref{sec:assembly}--\ref{sec:mono} and are not described here.

\subsection{The exact formula}
\label{sec:check_exact}
\begin{itemize}\itemsep2pt
\item \emph{Algebraic identities} (referee V1).  Lemma~\ref{lem:kin} was verified in a generic Lorentzian metric with no symmetry,
for generic $u$ and null $p$; Lemma~\ref{lem:fluid} for general $\rho,P$, with the converse that the identities hold iff $w=\tfrac13$
for $P=w\rho$, $\rho\propto T^4$; the off-shell identity behind Theorem~\ref{thm:exact} vanishes exactly on shell (V1-1--V1-3).
\item \emph{Exact solutions} (V1).  FRW and Milne: $\Tobs=T_o$ symbolically; Bianchi~I with a radiation test fluid (Kasner
$(\tfrac23,\tfrac23,-\tfrac13)$ and non-Kasner exponents), where the anisotropy is $O(1)$: $|\Tobs^{\rm exact}-T_oe^{-S}|\lesssim2\times10^{-31}$,
i.e.\ at quadrature precision (V1-4).  Linear Poisson gauge: the textbook Sachs--Wolfe $+$ Doppler $+$ ISW formula with observer terms is
recovered (V1-5).
\item \emph{Non-perturbative} (agent N1).  In the planar solver the photon energies are integrated from the geodesic equation and the
fluid temperatures are read off independently at emission and observation; $\max_{\nhat}|\Tobs/T_o-e^{-S}|=7.5\times10^{-15}$ at $\eps=0.01$ and
$1.4\times10^{-14}$ at $\eps=0.1$ ($\Theta_{\rm rms}=0.08$), and $\le7.5\times10^{-14}$ in every production run (N1-1).
\item \emph{Symbolic, through third order} (agent C2).  Route (i) (photon energy from $\hat\Gamma^0_{kk}$) and route (ii) ($e^{-S}$) agree
(a)~at the integrand level along an arbitrary perturbed ray at orders 1, 2, 3 modulo $\hat\nabla_\mu\hat T^{\mu\nu}=0$ only (C2-1--C2-3), and
(b)~at the level of the assembled, background-ray-expanded functionals, piece by piece (ray, obs, emi), orders 1, 2, 3 (C2-7).  The
two assembled forms in \code{L\_terms.pkl} agree numerically for the plane wave below to $4\times10^{-10}$ ($L_2$) and $9\times10^{-7}$ ($L_3$)
(A1-4).
\end{itemize}

\subsection{Einstein equations}
\label{sec:check_einstein}
\begin{itemize}\itemsep2pt
\item \emph{Internal} (C1).  Background parts vanish; the contracted Bianchi identity holds identically through $O(\eps^3)$, so the
constraints propagate; the first-order system reproduces the standard equations and the conservation laws hold on shell; the exact dilation
family solves the equations through $O(\eps^3)$ (C1-1--C1-6).
\item \emph{Independent confirmation} (referee V3).  A separately coded engine gives $E_{\mu\nu}$ term-by-term identical to C1's for all ten
components at orders 1, 2, 3 (V3-1; same method coded independently, compensated by the literature controls).  V3 confirmed the linear
solution, the operator structure at orders 2 and 3, all second-order kernels (agreement to $10^{-164}$ at random points), the algebraic
$\omega_2$ from relativistic Kelvin and its uniqueness, the superhorizon $\Phi_2$, the third-order contraction rules (brute-force Wick), the
contracted third-order sources, and $C_3$ re-derived independently at three random points to 30 digits (V3-2--V3-6).
\item \emph{Literature controls.}  Kohri--Terada \cite{KT}: $K_h=4f_{\rm RD}$, relative difference $5\times10^{-29}$ (V3-3, C1-9); Pitrou--Uzan--Bernardeau
\cite{PUB}: scalar sources (A.1)--(A.8) except the trace source (A.6), whose printed form is visibly garbled (our trace source is fixed by two
engines and by the exact dilation family), and the superhorizon initial data (3.12b,c) exactly (V3-4, C1-10); Lu et al.\ \cite{Lu}: the
radiation-era vector kernel is identical to $K_\omega$ (V3-5).
\item \emph{Closed forms} (agent C4).  The closed-form $\mathcal K_h$ equals Kohri--Terada (22) [with the modulus in the logarithm] to
$10^{-15}$ (25 digits with rational input) and their late-time form (25) exactly; the closed-form $\Phi_2$ agrees with 60-digit quadrature to
$3.5\times10^{-40}$; all second-order fields assembled from the closed forms satisfy C1's full second-order Einstein and conservation equations
with relative residual $5.5\times10^{-50}$ (C4-1--C4-3).
\item \emph{Plane-wave fields} (agent A1).  The orders-1--3 plane-wave fields generated from C1's engine satisfy C1's conservation equations to
$10^{-27}$/$5\times10^{-21}$ (orders 2/3) and reproduce N1's nonperturbative maps (Sec.~\ref{sec:check_endtoend}), which tests the equations
together with the initial-data prescription $\zeta^{(2)}_{\rm nl}(0)=\zeta^{(3)}_{\rm nl}(0)=0$ (A1-1).
\end{itemize}
A harmless erratum was found by V3: the small-$\eta$ helper for $\Phi_2$ in C1's \code{third\_order\_ic.py} uses $\tfrac3{10}S_\Phi(0)\eta^2$ instead
of $\tfrac1{10}$; $C_3$, the tadpoles and all other quoted results are independent of it.

\subsection{Dilation invariance}
\label{sec:check_dilation}
\begin{itemize}\itemsep2pt
\item \emph{The exact family.}  The rescaled FRW solution \eqref{eq:dilation} solves the field equations through $O(\eps^3)$ (C1-6) and
non-perturbatively (V3-7); it carries $\zeta_{\rm nl}=\ln B$ exactly and, with $B=e^{\eps c}$, is the $k\to0$ limit of the regular solution with
a free residual time rescaling at second order.
\item \emph{Symbolic cancellation in the observable} (C2).  For the family with $\lambda=1+\eps a_1+\tfrac{\eps^2}2a_2$, $a_1=-c/3$, $a_2$ free:
$\hat\sigma_{\mu\nu}\equiv0$ through third order; the assembled $L_1,L_2,L_3$ vanish in the shear \emph{and} the direct form, with
$\delta\eta_{e1}=c\eta_*/3$, $\delta\eta_{e2}=(-a_2+\tfrac29c^2)\eta_*$ inserted (the exact $\eta_*(1/\lambda-1)$ series); the clock gives
$\eta_o=\bar\eta_o/\lambda$ through second order for any $a_2$, hence $T_o=\bar T(\bar\eta_o)$ and $\Theta=0$ (C2-9).  In the $e^{-S}$
formulation the anisotropy part is automatic ($S$ involves only derivatives of $\hat u$), so the non-trivial content is the direct form and
the monopole at fixed proper time.  A factor-2 slip in C2's stored $\delta\eta_{e2}$ was caught by this check and corrected.
\item \emph{Squeezed limit of the initial data.}  $C_3(q\to0)=-\tfrac{32}9$, the value predicted by the dilation family (C1-14), re-derived by V3 (V3-6).
\item \emph{Zero-mode (tadpole) independence} (A1).  Regenerating the plane-wave fields with a different zero-mode convention
($\Psi_{2,0}=\Phi_{2,0}+c_g\eta$, a residual time reparametrisation) changes $\Phi_{2,0},\Psi_{2,0},\delta_{2,0}$ and all third-order fields at
$O(1)$ ($\Psi_{2,0}(8)$: $-1.21\to+4.79$) and $\delta\eta_{o2}$ ($4.411\to-11.589$), but every $a_{\ell0}$ of $\Theta^{(1,2,3)}$, the sky monopoles and $T_o$
are unchanged to $\le2\times10^{-14}$ (A1-3).  This is the numerical realisation of the statement that individual tadpoles are
convention dependent (V3 finds that the image-summed $C_3$ shifts by $8$ per unit shift of the tadpole $\langle\Phi_2\rangle$ along the allowed zero-mode family, V3-6) while the
observable, with clock and normalisation, is not.
\item \emph{First-order gradient mode.}  $\Delta_0=O(k^4)$, $\Delta_1=O(k^3)$ as $k\to0$ (C5-3): constant and gradient modes are invisible at first order.
\end{itemize}

\subsection{The decisive end-to-end test: plane wave versus the non-perturbative solver}
\label{sec:check_endtoend}
\paragraph{Set-up.}  For $\zeta=\eps\cos x$ (single plane wave, $k=1$) agent N1 solved the full planar Einstein--radiation system
non-perturbatively in comoving-orthogonal gauge (lapse fixed algebraically by Lemma~\ref{lem:fluid}), with regular non-linear initial data
at $t_i=0.01$, ray tracing to the hypersurface $\Tphys=T_*$, the observer clock at fixed proper time, and the normalisation of
Corollary~\ref{cor:theta}; $\Theta^{(1,2,3)}$ were extracted from 16 amplitudes $|\eps|\le0.1$ by even/odd polynomial fits, with an uncertainty
defined as the spread between fit subsets (N1-3).  Independently, agent A1 evaluated the analytic chain: C1's equations reduced to the plane
wave (fields to third order as exact Green's-function integrals, 40-digit arithmetic), C2's general-$\nhat$ shear-form lists $L_1,L_2,L_3$ with the
ray functionals \eqref{eq:k1def}--\eqref{eq:Dx2def}, the emission shifts \eqref{eq:deta_e}, the clock and displacement of Sec.~\ref{sec:obsevent}
(including the $O(t_i^2)$ initial-event terms) and the bookkeeping \eqref{eq:theta_orders} including $\ell=0,1$.  No fitting of A1 to N1 is
involved.  Configurations: $\bar\eta_o/\eta_*=4$ ($\bar\eta_o=8$) and $2$ ($\bar\eta_o=4$), $\eta_*=2$; observer at $x_o=0$ (maximum of $\zeta$; even
$\ell$ only) and $x_o=\pi/2$ (generic; odd $\ell$ and dipole).  Multipoles $a_{\ell0}=\sqrt{4\pi(2\ell+1)}\,\tfrac12\int\Theta P_\ell\,\dd\mu$ about the
wave vector.

\begin{table}[t]
\centering\small
\begin{tabular}{@{}lrr rrr@{}}
\hline
configuration & order & $\ell$ & A1 (analytic chain) & N1 (non-perturbative) & A1$-$N1\\ \hline
$\bar\eta_o=8$, $x_o=0$ & 1 & 2 & $-2.1847203553$ & $-2.1847203553$ & $-4.4\times10^{-12}$\\
 & 1 & 4 & $+1.2552456282$ & $+1.2552456282$ & $-3.6\times10^{-12}$\\
 & 1 & 6 & $-4.7973477367\times10^{-2}$ & $-4.7973477364\times10^{-2}$ & $-3.4\times10^{-12}$\\
 & 2 & 2 & $-4.2224942186$ & $-4.2224942193$ & $+6.5\times10^{-10}$\\
 & 2 & 4 & $+1.3785439774$ & $+1.3785439781$ & $-7.4\times10^{-10}$\\
 & 2 & 6 & $-1.2595382846$ & $-1.2595382817$ & $-2.9\times10^{-9}$\\
 & 3 & 2 & $-1.2889987902\times10^{1}$ & $-1.2889987968\times10^{1}$ & $+6.6\times10^{-8}$\\
 & 3 & 4 & $-4.7070907021\times10^{1}$ & $-4.7070907098\times10^{1}$ & $+7.7\times10^{-8}$\\
 & 3 & 6 & $+6.2905106702$ & $+6.2905105363$ & $+1.3\times10^{-7}$\\ \hline
$\bar\eta_o=4$, $x_o=0$ & 1 & 2 & $+2.0752530981\times10^{-1}$ & $+2.0752530981\times10^{-1}$ & $+4.3\times10^{-12}$\\
 & 2 & 2 & $-1.3572165380$ & $-1.3572165385$ & $+5.4\times10^{-10}$\\
 & 2 & 4 & $-2.2995087658\times10^{-1}$ & $-2.2995087720\times10^{-1}$ & $+6.1\times10^{-10}$\\
 & 3 & 2 & $-4.8330628410$ & $-4.8330627544$ & $-8.7\times10^{-8}$\\
 & 3 & 4 & $-3.0635108209$ & $-3.0635106964$ & $-1.2\times10^{-7}$\\ \hline
$\bar\eta_o=8$, $x_o=\pi/2$ & 1 & 1 & $+2.0812147132$ & $+2.0812147132$ & $+1.1\times10^{-11}$\\
 & 1 & 3 & $+2.0917805127$ & $+2.0917805128$ & $-5.1\times10^{-11}$\\
 & 2 & 2 & $+3.0669017921$ & $+3.0669017899$ & $+2.2\times10^{-9}$\\
 & 2 & 4 & $+6.5664035508$ & $+6.5664035492$ & $+1.6\times10^{-9}$\\
 & 3 & 1 & $-5.1723421459\times10^{1}$ & $-5.1723420730\times10^{1}$ & $-7.3\times10^{-7}$\\
 & 3 & 3 & $-4.3006733540\times10^{1}$ & $-4.3006736906\times10^{1}$ & $+3.4\times10^{-6}$\\ \hline
$\bar\eta_o=4$, $x_o=\pi/2$ & 1 & 1 & $-2.3149107586\times10^{-1}$ & $-2.3149107587\times10^{-1}$ & $+3.8\times10^{-12}$\\
 & 2 & 2 & $-2.1076209609$ & $-2.1076209611$ & $+2.1\times10^{-10}$\\
 & 3 & 1 & $+1.3059761781\times10^{1}$ & $+1.3059761844\times10^{1}$ & $-6.4\times10^{-8}$\\
 & 3 & 3 & $-3.6965575489\times10^{-1}$ & $-3.6965566762\times10^{-1}$ & $-8.7\times10^{-8}$\\ \hline
\end{tabular}
\caption{Plane-wave comparison of the analytic chain (agent A1) with the non-perturbative planar solver (agent N1): $a_{\ell0}$ of
$\Theta^{(r)}$, $r=1,2,3$ ($\ell=1$ is the dipole, kept here because the comparison is made before dipole removal).  Full lists
($\ell\le12$): \code{code/A1\_planewave/compare\_*.log}.}
\label{tab:a1n1}
\end{table}

\begin{table}[t]
\centering\footnotesize\setlength{\tabcolsep}{4pt}
\begin{tabular}{@{}lccc@{}}
\hline
configuration & $\Theta^{(1)}$ & $\Theta^{(2)}$ & $\Theta^{(3)}$\\ \hline
$\bar\eta_o=8$, $x_o=0$ & $1.5\times10^{-10}$ ($3.6\times10^{-9}$) & $1.3\times10^{-7}$ ($3.2\times10^{-8}$) & $4.0\times10^{-6}$ ($2.1\times10^{-4}$)\\
$\bar\eta_o=4$, $x_o=0$ & $1.3\times10^{-11}$ ($6.8\times10^{-12}$) & $4.5\times10^{-9}$ ($9.7\times10^{-10}$) & $2.7\times10^{-7}$ ($4.0\times10^{-7}$)\\
$\bar\eta_o=8$, $x_o=\pi/2$ & $2.7\times10^{-8}$ ($5.3\times10^{-7}$) & $4.0\times10^{-7}$ ($9.2\times10^{-6}$) & $1.8\times10^{-3}$ ($3.1\times10^{-2}$)\\
$\bar\eta_o=4$, $x_o=\pi/2$ & $2.1\times10^{-11}$ ($6.7\times10^{-10}$) & $2.4\times10^{-9}$ ($5.4\times10^{-9}$) & $3.1\times10^{-7}$ ($3.9\times10^{-5}$)\\ \hline
map rms ($\bar\eta_o=8$, $x_o=0$) & $0.71$ & $1.31$ & $14.2$\\ \hline
\end{tabular}
\caption{Maximum over the sky of $|\Theta^{(r)}_{\rm A1}-\Theta^{(r)}_{\rm N1}|$, with N1's extraction uncertainty in parentheses.}
\label{tab:a1n1maps}
\end{table}

\paragraph{Result.}  Tables~\ref{tab:a1n1} and \ref{tab:a1n1maps}: the analytic chain reproduces the non-perturbative maps at all three
orders in all four configurations, within N1's uncertainty: first order to $\sim10^{-11}$, second order to $\sim10^{-9}$ in $a_{\ell0}$, third order
to $\sim10^{-7}$ ($x_o=0$) and $\sim10^{-6}$ ($x_o=\pi/2$, where N1's own uncertainty is $3\times10^{-2}$) (A1-2).  The second-order map differences for
$x_o=0$ exceed N1's quoted spread by $\sim4$ but sit at N1's stated noise floor; a test that uses no N1 fitting at all settles this:
the residual $R(\eps)=\Theta_{\rm N1}(\eps)-[\eps\Theta^{(1)}+\tfrac{\eps^2}2\Theta^{(2)}+\tfrac{\eps^3}6\Theta^{(3)}]_{\rm A1}$ scales cleanly as $\eps^4$ for all 16
amplitudes ($\max|R|/\eps^4=11.6\ldots14.5$ smoothly over $-0.1\le\eps\le0.1$; $|R|=1.2\times10^{-7}$ at $\eps=0.01$), i.e.\ no lower-order
remainder (A1-6).  The gauge-invariant observer quantities agree as well: $T_o$ and $\sky{\Tobs}$ through $O(\eps^2)$, e.g.\ for
$\bar\eta_o=8$, $x_o=0$: $T_o=0.125-0.0550223132\,\eps+0.0041441746\,\eps^2$ (A1) vs $0.125-0.0550222904\,\eps+0.0041441914\,\eps^2$ (N1 fit), with N1 fit
uncertainties $10^{-8}$--$5\times10^{-7}$.  (The coordinate time of observation is gauge dependent and is not compared.)
A1's own numerics are converged (Chebyshev $N=48/64/96$ identical to all printed digits).

\paragraph{What it tests.}  Every link of the chain enters: C1's second- and third-order equations and the initial-data prescription,
the tadpoles at third order, C2's $L_1,L_2,L_3$ including lensing, time delay, aberration, the direction map with the canonical triad, the
emission surface to second order, the clock and observer displacement to second order, and the normalisation \eqref{eq:theta_orders}.
On the way two errors were found and corrected, neither in C1's equations nor in C2's term lists: a zero-mode weight in A1's own
conversion $\delta\to\delta_T$ (detected by C2's conservation rules and by $T_o$ at $O(\eps^2)$, A1-7) and the documentation sign slip of
Remark~\ref{rem:signslip}, which with the printed $+$ sign spoils the null condition at $O(1)$ (A1-4).  The comparison also fixed the observer
label convention (Sec.~\ref{sec:obsevent}): with the label taken at $\eta_{\rm init}$ instead of $\eta\to0$ the observer is a different fluid element,
which moves $\Theta^{(2)}$'s $a_{20}$ by $7\times10^{-5}$, beyond N1's uncertainty (A1-5).

\subsection{First-order controls}
\label{sec:check_first}
Agent C5 (Sec.~\ref{sec:linear}): the transfer functions satisfy the first-order equations and constraints exactly; the Rayleigh rules in the
convention of Sec.~\ref{sec:dirconv} were checked by sphere quadrature ($3\times10^{-14}$, $\ell\le6$), including the observer monopole ($\ell=0$ only) and
observer Doppler ($\ell=1$ only); the prefactor of $C^{(11)}_\ell$ and the observer and ISW terms by an independent real-space computation of the
correlation function ($6\times10^{-15}$); the ISW closed form against quadrature ($2\times10^{-18}$), its residue formulas for the Ci/Si/Cin
coefficients ($<10^{-169}$) and the vanishing of the $\ln w$ coefficient (regularity at the observer) symbolically; the limits
$k\eta_*\to0$ (Sachs--Wolfe $-\tfrac13$, plateau $A_s/9$), $k\to\infty$ and $k\to0$; and $\langle\sky{m_1}a_{\ell m}\rangle=0$ ($\ell\ge1$),
$\langle\sky{m_1}^2\rangle=C^{(11)}_0/4\pi$ (C5-1--C5-4).  The first-order map is in addition tested end to end by the plane-wave comparison
($\sim10^{-11}$).

\subsection{Harmonic machinery controls}
\label{sec:check_harmonic}
Agent C3 (Sec.~\ref{sec:harmonic}): rules (A-$s$), (A-$\lambda$), (A-$L$), the type-C $3j$ sums and the reduction (E1) against direct quadrature
($10^{-14}$--$10^{-15}$); (B1) with both pairings against an exact real-space Wick computation ($2.6\times10^{-12}$) and by Monte Carlo on the
sphere; the type-D Fourier/angular-gradient identity ($3\times10^{-12}$); the prefactors 2 and 3 of \eqref{eq:loops}.  The decisive control is
the lensed $C_\ell$ (Sec.~\ref{sec:lensing_control}): the bookkeeping \eqref{eq:cl_target} reproduces Hu's $(22)$ and $(13)$ lensing terms, extended to
correlated $\phi,\Theta$, algebraically ($2.5\times10^{-16}$) and by Monte Carlo of the remapped sky ($2$--$5\%$, statistical) (C3-1--C3-9).
Formulas (E2), (E3) and $T^{\rm(ii)}$ are derived by Wick counting consistent with the verified ones but have not been tested numerically (C3-10).

\subsection{What the checks do not cover}
\label{sec:check_gaps}
\begin{itemize}\itemsep2pt
\item A single plane wave probes only the axisymmetric ($m=0$ about $\vk$) part of the kernels $K_r$, and only the collinear configurations
$\vk_i\parallel\vk$ of $K_2$, $K_3$ (with the $x_o=0$ configurations restricted to even $\ell$); the planar solver contains no vector or tensor
modes beyond a homogeneous $h$, so the $\omega_2$ terms and the non-collinear $h_2$ terms of the map are untested end to end.
\item Only $\bar\eta_o/\eta_*=2$ and $4$ were compared, far from a realistic $\bar\eta_o/\eta_*$; the third-order comparison at ratio 4 and
$x_o=\pi/2$ is limited by N1's extraction uncertainty ($3\times10^{-2}$ absolute on a map of rms 24).
\item The loop integrals over $\vk_1,\vk_2$ and $\vq$, the harmonic transforms of the actual second- and third-order kernels, and the
monochromatic evaluation (including the delicate limits listed in Sec.~\ref{sec:monodef}) are so far checked only by internal consistency;
their checks are recorded in Secs.~\ref{sec:assembly}--\ref{sec:mono}.  The squeezed-limit predictions of agent V2 (adiabatic-mode
transformation, App.~\ref{app:ledger}, V2-1) await comparison with the assembled $K_2$.
\item The tetrad-rotation invariance of $\tfrac14C^{(22)}+\tfrac13\Re C^{(13)}$ (V1-6) is verified only for an exact $\ell=2$ example; it is not needed if
the canonical triad is used consistently, as here.
\item The literature agreement for the scalar second-order trace source rests on our two engines and the dilation family only, since
the printed PUB (A.6) is garbled.
\end{itemize}
